\section{Future work}
\label{sec:future-work}

TODO sketch general net of equivalences in Proof pearls aswell (image)

The two results of this paper are intended as the first pieces of a
larger, ongoing effort within CSLib to build a formalized foundation
for research on untyped lambda calculus. We outline the planned
directions below.

\subsection{Further equivalence results}
\label{subsec:further-equivalences}

Crole-style results, showing several natural but superficially
different definitions of the same notion to coincide, are not unique
to alpha equivalence. We plan analogous formalizations for SK
combinator calculus and other binding representations (locally
nameless terms, nominal terms), each intended to slot into CSLib
alongside the results of this paper and to be connected to them by
explicit isomorphisms in the style of Section~\ref{sec:de-bruijn}.

\subsection{A formalized Barendregt}
\label{subsec:formalized-barendregt}

Barendregt's \emph{The Lambda Calculus: Its Syntax and
Semantics}~\cite{barendregt1985a} remains the standard reference
for the untyped lambda calculus, and a substantial fraction of its
content --- confluence, standardization, the theory of reduction,
combinatory completeness, and the associated notions of models ---
has never been formalized in a single, coherent development. We see
CSLib's binding infrastructure as a natural basis for such an effort,
which we tentatively refer to as a "formalized Barendregt". Beyond
its intrinsic value as a reusable library, a sufficiently complete
formalized Barendregt would put open or only recently resolved
conjectures discussed in Barendregt's later survey work within reach
of machine-checked treatment.

\subsection{Maintenance and reuse}
\label{subsec:maintenance}

Throughout, the intent is for this work to be maintained as part of
CSLib rather than as a standalone, one-off development: subject to
code review, kept up to date with CSLib and Mathlib, and structured
so that later projects can depend on it directly.
